% ============================================================
% V. DISCUSSION and VI. CONCLUSION (condensed to the IEEE conference length)
% Every number is taken from docs/RESULTS.md and the committed results files.
% ============================================================

\section{Discussion}
\label{sec:discussion}

\emph{Why the proposed model does not win.} On the stream TFT\,+\,EWC never beat linear extrapolation: fade over one to three checks is smooth, the groups are few and small (21--27 and 135--225 training samples per task in S1 and S6), and the TFT barely uses its history. EWC helps where there is something to protect (negligible forgetting across the near-identical S1 cells, substantial across the S6 conditions), most clearly on the residual target, where it cuts forgetting over six horizons by about two thirds; a shift-adaptive penalty added nothing, and at horizons 1--3 a BiLSTM or Transformer with EWC is as accurate as the residual TFT. On unseen conditions the continual constraint itself is the main cost (joint training is 0.32--0.44~pp more accurate): with data this small, storing every earlier group is cheap, so the constraint that motivates continual learning did not bind.

\emph{What matters.} Building on the best simple forecaster mattered most, and the ablation points to inputs, not architecture: diagnostic features hurt the TFT and the calendar time between checks helped, consistent with fade per cycle rising with that time at 50~$^\circ$C in all 20 S6 cells. For intervals, the calibration rows mattered as much as the model: residuals from mid-life validation origins of seen conditions understate the error on a new one, residuals from held-out conditions reproduce it. With few cells per condition, 21 dense ageing-cycle inputs added more variance than information, and first-life cells of other chemistries did not transfer to the LFP cells of S6.

\emph{Reproducibility and reuse.} As significance results changed across hardware (Section~\ref{sec:res-repro}), we recommend recording the environment, tuning with enough runs to separate neighbouring grid values, and requiring replicated differences above a practical threshold. The reuse classes are informative only when the horizon crosses class boundaries (S1 stays in C0, the S6 test periods lie in C3), and for a population unlike any seen, such as the 100\% depth-of-discharge condition, neither forecasts nor intervals can be trusted, which pooled coverage hides. The studies of Section~II estimate current SOH or SOC on other cells and SOH definitions, so their errors are not directly comparable with ours.

\emph{Limitations.} The task was changed to forecasting after the first dataset was inspected; the second-life datasets are small; S6 is aged at 50~$^\circ$C and its ``second life'' continues the same protocol; SOH is relative to each cell's first valid measurement; forecasts are conditional on planned usage; and the ablation and every later experiment were decided after the main results (all except the ablation were committed before they were run). The single-fit regressors see all tasks at once, and the direct literature backbones were run for horizons 1--3 only. The TFT hyperparameters used on unseen conditions were tuned with validation rows of all seven S6 conditions, which come from one laboratory, with only training randomness varied. Leave-one-condition-out intervals were run for gradient boosting and extrapolation only.

\section{Conclusion}
\label{sec:conclusion}

We evaluated a TFT with EWC as a continual learner for forecasting the SOH of second-life cells, on a leakage-controlled benchmark with pre-registered experiments and replicated differences. Learning a shrunk correction to linear extrapolation made it as accurate as extrapolation on the stream and more accurate on unseen usage conditions, where gradient boosting and the same TFT trained jointly were clearly better: at this data scale, learning continually cost more than it gave. Conformal intervals held on unseen conditions only when calibrated on conditions held out from fitting, and not outside the training range; refitting and tracking cut the interval score by 29\%, and 87--93\% of forecasts fell within one reuse class, over 99\% of them correctly. Future work should test continual learning where earlier data cannot be stored and populations are larger, and calibrate intervals across laboratories and chemistries.
